\documentclass[hyperref={pdfpagelabels=false},table, handout]{beamer}

\input{lec_common}
\begin{document}

\part{Unit 09: Tests and  Profiling}



\begin{frame}{Why Tests?}

\begin{center}
        {\Huge ?}
\vfill\pause
\end{center}
\begin{itemize}
        \item Because you do them \alert{anyway}
        \item Because it will keep your code \alert{alive}.
\end{itemize}
\end{frame}


\begin{frame}[fragile]{Automated Tests}
Automated tests 
\begin{itemize}
\item ensure a constant (high) quality standard of your code
\item serve as a documentation of the use of your code
\item document the test cases $\rightarrow$ test protocol
\end{itemize}
\end{frame}

\begin{frame}[fragile]{Example}
Let's assume we want to test an implementation of the bisection algorithm:
\begin{python}
def bisect(f, a, b, tol=1.e-8):
    if f(a) * f(b) >0:
       raise ValueError("Incorrect initial [a, b]")   
    for i in range(100):
        c = (a + b) / 2.
        if f(a) * f(c) <= 0:
           b = c
        else:
           a = c
        if abs(a - b) < tol:
           return (a + b) / 2
    raise Exception('No root found')   
\end{python}

\end{frame}


\begin{frame}[fragile]{Example (Cont.)}
We check a problem with a known solution.\\
Does the code find a zero of $f(x)=x$? 
\begin{python}
def test_identity():
    result = bisect(lambda x: x, -1., 1.) 
    expected = 0.
    assert allclose(result, expected),'expected zero not found'

test_identity()
\end{python}
Note the command \pyth{allclose}
\end{frame}
\begin{frame}[fragile]{Example (Cont.)}
Does the code handle wrong input correctly? \vfill
\begin{python}
def test_badinput():
  try:
    bisect(lambda x: x,0.5,1)
  except ValueError:
    pass
  else:
    raise AssertionError()

test_badinput()
\end{python}


\end{frame}


\subsection{unittests}
\begin{frame}[fragile]{Unittest Test Cases}
You will want to \alert{put your tests together} and automatize them:
\begin{python}
from bisection import bisect
import unittest

class TestIdentity(unittest.TestCase):
    def test(self):
        result = bisect(lambda x: x, -1.2, 1.,tol=1.e-8)
        expected = 0.
        self.assertAlmostEqual(result, expected)

if __name__=='__main__':
    unittest.main()
\end{python}
\end{frame}

\begin{frame}[fragile]{Test results}
Here the results with two different tolerance parameters:

\tiny
\begin{lstlisting}
Ran 1 test in 0.002s

OK
\end{lstlisting}
\normalsize

... and here with a loose tolerance
\tiny
\begin{lstlisting}
F
======================================================================
FAIL: test (__main__.TestIdentity)
----------------------------------------------------------------------
Traceback (most recent call last):
  File "<ipython-input-11-e44778304d6f>", line 5, in test
    self.assertAlmostEqual(result, expected)
AssertionError: 0.00017089843750002018 != 0.0 within 7 places

----------------------------------------------------------------------
Ran 1 test in 0.004s

FAILED (failures=1)

\end{lstlisting}
\normalsize
\end{frame}

\begin{frame}[fragile]{Unittest: Grouping tests together}
We recommend to group tests together
\begin{python}
class TestIdentity(unittest.TestCase):
     def identity_fcn(self,x):
         return x
     def test_functionality(self):
         result = bisect(self.identity_fcn, -1.2, 1.,tol=1.e-8)
         expected = 0.
         self.assertAlmostEqual(result, expected)
     def test_reverse_boundaries(self):
         result = bisect(self.identity_fcn, 1., -1.)
         expected = 0.
         self.assertAlmostEqual(result, expected)
     def test_exceeded_tolerance(self):
         tol=1.e-80
         self.assertRaises(Exception, bisect, self.identity_fcn, 
                             -1.2, 1.,tol)        
if __name__=='__main__':
       unittest.main()
\end{python}
\end{frame}
\begin{frame}[fragile]{unittest.TestCase.assertRaises}
Note the method \pyth{assertRaises}:\vfill
\begin{itemize}
\item Exception: the expected exception type
\item bisect: the function to be called
\item self.identity\_fcn, -1.2, 1.,tol: the parameters of this function
\end{itemize}
\vfill
\end{frame}
\begin{frame}[fragile]{Unittest: Preparing tests}
Sometimes the execution of tests need some preparation. See special method \pyth{setUp}:
\begin{python}
class TestFindInFile(unittest.TestCase):
     def setUp(self):
         file = open('test_file.txt', 'w')
         file.write('aha')
         file.close()
         self.file = open('test_file.txt', 'r')
     def tearDown(self):
         os.remove(self.file.name)
     def test_exists(self):
         line_no=find_string(self.file, 'aha')
         self.assertEqual(line_no, 0)
     def test_not_exists(self):
         self.assertRaises(NotFoundError, 
                 find_string,self.file, 'bha')

if __name__=='__main__':
                unittest.main() 
\end{python}        

\end{frame}

\begin{frame}[fragile]{Unittest: Preparing tests\hfill(Cont.)}
... and here the functions which we want to test
\begin{python}
class NotFoundError(Exception):
      pass

def find_string(file, string):
    for i,lines in enumerate(file.readlines()):
       if string in lines:
          return i
    raise NotFoundError('String {} not found in File {}'.
                               format(string,file.name)) 
\end{python}        

\end{frame}



\section{Profiling}


\begin{frame}[fragile]{Slow code}
The code is slow. Why?

\vfill

Three possible reasons:
\begin{itemize}
\item The problem is big \hfill \emph{Not much to do about this}
\item The algorithm is inherently slow \hfill \emph{Nor this}
\item The implementation is bad (recursive, etc.) \hfill \emph{We can fix this!}
\end{itemize}
\end{frame}

\begin{frame}[fragile]{Finding the slow parts}
Different parts of the code are differently fast.

\vfill

\emph{Profiling} means that we run the program and measure how much time is spent in each function.

\vfill

A good way to do this in Python is to use the \pyth!cProfile! module.

\end{frame}

\begin{frame}[fragile]{The timeit module}
The \pyth!timeit! module is ideal for measuring execution time for short statements:

\vfill
 What is faster?
\begin{python}
def var1():
   [i**2 for i in range(10000) if i%2==0]
     
def var2():
   li=[]
   for i in range(10000):
      if i%2==0:
          li.append(i**2)
\end{python}

When this code is saved in \pyth!timeit_ex.py!, then

\begin{python}
import timeit as ti
t1=ti.Timer('te.var1()','import timeit_ex as te')
t1.timeit(100) # a hundred executions
t2=ti.Timer('te.var2()','import timeit_ex as te')
t2.timeit(100) # a hundred executions
\end{python}

gives us the answer....

\vfill

\end{frame}

\begin{frame}[fragile]{The timeit module \hfill (Cont.)}
 The \pyth!timeit! module provides a \pyth!Timer! class
 with the signature
\begin{python}
class Timer(stmt[,setup]):  
\end{python}
\vfill
where 

\pyth!stmt! is a string containing the statement(s) to be executed.\\
\pyth!setup! is a string with the statement(s) executed beforehand, i.e. not affecting the time measurement.

\vfill

It has a method \pyth! timeit! which takes the numbers of repetitions as argument
\vfill
\end{frame}


\begin{frame}{Other timeing Methods}
\vfill
For other timeing methods consult the course book.
\vfill
\end{frame}
\end{document}
